\documentclass[11pt,twoside,openright]{book}

\usepackage[T1]{fontenc}
\usepackage[utf8]{inputenc}
\usepackage{amsmath,amssymb,amsthm,mathtools,bm}
\usepackage{mathpazo}
\usepackage{microtype}
\usepackage[inner=1.35in,outer=1.1in,top=1.25in,bottom=1.35in]{geometry}
\usepackage{booktabs}
\usepackage{fancyhdr}
\usepackage[hidelinks]{hyperref}
\usepackage{cleveref}

\setlength{\parindent}{1.4em}
\setlength{\parskip}{0pt}
\linespread{1.07}
\emergencystretch=3em

\pagestyle{fancy}
\fancyhf{}
\fancyhead[LE]{\small\thepage\quad\textsc{A Cinema of Individual Editions}}
\fancyhead[RO]{\small\nouppercase{\leftmark}\quad\thepage}
\renewcommand{\headrulewidth}{0.3pt}
\fancypagestyle{plain}{\fancyhf{}\fancyfoot[C]{\small\thepage}\renewcommand{\headrulewidth}{0pt}}

\theoremstyle{plain}
\newtheorem{theorem}{Theorem}[chapter]
\newtheorem{proposition}[theorem]{Proposition}
\newtheorem{lemma}[theorem]{Lemma}
\newtheorem{corollary}[theorem]{Corollary}
\theoremstyle{definition}
\newtheorem{definition}[theorem]{Definition}
\newtheorem{example}[theorem]{Example}
\theoremstyle{remark}
\newtheorem{remark}[theorem]{Remark}

\newcommand{\Capac}{\mathrm{Cap}}
\newcommand{\Rob}{\mathrm{Rob}}
\newcommand{\Inv}{\mathrm{Inv}}
\DeclareMathOperator{\loss}{loss}

\title{\Huge\bfseries A Cinema of Individual Editions\\[1.2ex]
\large Chapter Thirteen: The Lifetime Repertoire\\[0.4ex]
\large Working Draft}
\author{Flyxion\\[0.3ex]\small Independent Researcher}
\date{November 2026}

\begin{document}

\frontmatter
\maketitle
\tableofcontents
\mainmatter
\setcounter{chapter}{12}

\chapter{The Lifetime Repertoire}
\label{ch:repertoire}

A viewer has a finite viewing life. The ledger of Chapter~12 will record, for one person, only as many editions as that life has room for, and the question this chapter asks is what the room is, how it is best spent among films, revisits, and release dates, and what a production must guarantee about the editions that fill it. The answer turns on a trade that the book has been circling since its first chapter. A person who wants never to see the same film twice can do so by seeing a great many different films, or by seeing a smaller number of films in editions that differ enough to count as different. The first costs productions. The second costs the tolerance of Chapter~2 and the capacity of Chapter~9.

The chapter works with the lawn scene and with the measured length of the passage of Chapter~6, an 87-minute film of $5{,}222$ seconds. The viewing rates are assumptions chosen to be easy to change, and a person can have any rate. The space of possible films is never enumerated, and the chapter shows that it does not need to be.

\section{The Budget}\label{sec:budget13}

Let a viewer watch $w$ hours per week for $y$ years and let a viewing last $h$ hours. The number of viewings is $N=\lfloor 52wy/h\rfloor$. With $y=60$ and $h=1.4506$, which is the length of the measured film, \Cref{tab:budget13} gives the budget at three rates.

\begin{table}[ht]
\centering
\begin{tabular}{@{}rrrr@{}}
\toprule
Hours per week & Lifetime hours & Viewings $N$ & Viewings per year\\
\midrule
2 & 6{,}240 & 4{,}301 & 71.7\\
6 & 18{,}720 & 12{,}905 & 215.1\\
14 & 43{,}680 & 30{,}112 & 501.9\\
\bottomrule
\end{tabular}
\caption{Lifetime viewings of an 87-minute film over sixty years at three weekly rates.}\label{tab:budget13}
\end{table}

All later counts use the first row, $N=4{,}301$. The budget is a bound on viewings, and it says nothing yet about which films fill it. A viewer who never revisits and wants no repeat of a film needs $4{,}301$ different films, which is about seventy-two new films a year for sixty years.

\section{Revisits and the Novelty Horizon}\label{sec:horizon13}

A revisit is a viewing of a film already seen. Whether it is a repeat depends on the edition. If the edition shown is not perceptually equivalent to any edition the viewer has seen of that film, then by the standard of Chapter~2 the viewer has seen something new, though the story is the same. The \emph{novelty horizon} of a production at tolerance $\varepsilon$ is the number of viewings a single viewer can have of it with every edition pairwise more than $\varepsilon$ apart. By the definition of the packing number in Chapter~2, it is $M_\varepsilon$. Two uses of viewing time both go by the name of revisiting, and the count should say which it is using. Chapter~1 budgets time for seeing an earlier film again, which is a repeat of the same work. This chapter counts a perceptibly different edition of the same work toward the novelty horizon, which is a new viewing of the work and not a repeat of what was seen. Both are legitimate uses of the budget.

\begin{proposition}[Lifetime catalogue]\label{prop:life}
Suppose every production in a catalogue has novelty horizon at least $M$ at the viewer's tolerance. A viewer with $N$ viewings can avoid any perceptual repeat within a production with $\lceil N/M\rceil$ productions, and cannot do so with fewer. At a constant viewing rate of $N/y$ per year, the catalogue must grow at no fewer than $N/(yM)$ new productions a year to keep the viewer supplied.
\end{proposition}

\begin{proof}
Each production supplies at most $M$ viewings without a repeat, so $\lceil N/M\rceil$ productions are necessary, and they suffice because each can supply $M$.
\end{proof}

\Cref{tab:horizon13} gives the horizons of production A and production B after the agreement rule, from the packing numbers of Chapter~2, and the catalogue that a person watching two hours a week would need. The first row of each production is exact inequivalence, and the others are tolerances of one and two slots.

\begin{table}[ht]
\centering
\begin{tabular}{@{}llrrp{2.6cm}p{1.6cm}@{}}
\toprule
Prod. & $\varepsilon$ & Horizon $M_\varepsilon$ & Weeks to exhaust & Productions for $4{,}301$ viewings & New films a year\\
\midrule
A & 0 & 238 & 172.6 & 19 & 0.30\\
A & 1 & 54 & 39.2 & 80 & 1.33\\
A & 2 & 16 & 11.6 & 269 & 4.48\\
B & 0 & 96 & 69.6 & 45 & 0.75\\
B & 1 & 32 & 23.2 & 135 & 2.24\\
B & 2 & 12 & 8.7 & 359 & 5.97\\
\bottomrule
\end{tabular}
\caption{Novelty horizon and the catalogue a lifetime of $4{,}301$ viewings needs, at $1.38$ viewings a week. Without revisits the same viewer needs $4{,}301$ films.}\label{tab:horizon13}
\end{table}

The compression is the ratio of the last column of viewings to the number of productions, and it equals the horizon. A viewer who accepts a tolerance of one slot in production A can fill sixty years with eighty productions in place of $4{,}301$ films, and the production has to release one new film every nine months in place of one every five days. At tolerance two the horizon of A is sixteen, and a viewer who watches weekly exhausts a production's perceptually distinct editions in a little under twelve weeks.

This is the arithmetic that the framing of Chapter~1 asked for, and it comes with two qualifications that the table does not carry. A revisit is not the first viewing of a new story, and the proposition counts only the perceptual novelty of the edition. A viewer who remembers the film is not seeing the film for the first time at any tolerance, and Chapter~14 takes this up as the dependence of meaning on order. And the horizons are those of the toy scene, with a Hamming dissimilarity standing in for a measured one.

\section{Selecting Without Enumerating}\label{sec:select13}

A repertoire is drawn from the editions of many productions, and for a real feature the class tuples cannot be listed. A feature with a thousand slots and twenty classes at each has a space that no enumeration could touch. Counting and selection do not need one.

The admissible class tuples under the direction rule form the paths through a layered graph, and the number of paths from each class at each slot to the end of the scene is computed in one backward pass: the number from a class at slot $j$ is the sum of the numbers from the classes at slot $j+1$ that its direction accepts.

\begin{proposition}[Unranking]\label{prop:unrank}
Let the admissibility of a class tuple depend only on adjacent pairs, as the direction rule does. The number of admissible tuples is computed with $O(mq^2)$ additions of integers, for $q$ classes at each of $m$ slots, and the $r$th admissible tuple in lexicographic order is produced from $r$ with $O(mq)$ comparisons. Every admissible tuple is produced for exactly one $r$.
\end{proposition}

\begin{proof}
The table of counts decomposes the set of tuples by the class chosen at each slot in turn. At slot $j$, subtracting the counts of the classes passed over from $r$ identifies the class whose block of tuples contains the $r$th, and the remainder indexes within the block. The map is a bijection between the integers below the total and the admissible tuples, since the blocks partition the set.
\end{proof}

The procedure was checked against full enumeration for all six pools of the lawn scene, the three productions at the premiere and after the agreement rule, and for every $r$ it returns the same tuple as the enumeration: counts of $494$, $238$, $320$, $96$, $2$ and $1$. It was then run without enumeration on a synthetic scene of forty slots and ten classes, with a count of thirty-four digits, and on one of a thousand slots and twenty classes, whose count has $1{,}147$ digits. Twenty ranks drawn at random in the larger scene were unranked into twenty distinct valid tuples in $0.16$ seconds. The rules of Chapter~5 that reach past neighbors enlarge the table by the width of the frontier, which was one, two and three in the lawn scene, and the same argument applies.

A repertoire can therefore be a list of ranks. Selecting an edition is choosing an integer below the count, which a viewer's receipt can hold in the compact encoding of Chapter~12.

\section{Scheduling a Repertoire}\label{sec:schedule13}

Selection has to meet two demands at once. At each moment the editions in the hands of the audience must be distinct, which is the demand of the earlier chapters. And over time each viewer's own sequence must avoid repeats for as long as the pool allows. A cyclic schedule over the ranks meets both.

\begin{proposition}[Cyclic schedule]\label{prop:cyclic}
Let the admissible class tuples of a production be ranked $0,\dots,K-1$, let $U\le K$ viewers have distinct offsets $a_v$ in the integers modulo $K$, and let $s$ be coprime to $K$. If viewer $v$ is shown rank $(a_v+ks)\bmod K$ at the $k$th round, then at every round the $U$ viewers hold pairwise distinct editions, and every viewer sees each of the $K$ editions exactly once in $K$ rounds.
\end{proposition}

\begin{proof}
At a fixed round the map from offset to rank is a translation of the integers modulo $K$, which is injective. For a fixed viewer the map from round to rank is a multiplication by a unit followed by a translation, which is a bijection on the $K$ rounds.
\end{proof}

The proposition guarantees exact inequivalence only, and the perceptual quality depends on the stride. In ranked order, consecutive tuples are neighbors in the last slot. With stride one, a viewer's consecutive editions differ in $1.45$ slots on average in A and $1.91$ in B, and sometimes in one. A search over the strides coprime to $K$ finds, for the $238$ editions of A, stride $75$ with consecutive editions differing in at least three slots and $4.18$ on average, and for the $96$ of B stride $47$ with every consecutive pair differing in all five slots. These are distances between successive viewings and not guarantees that every pair of a viewer's editions is far apart, which would be the packing of the previous section and a shorter horizon.

The schedule also fixes what a tolerance costs for a simultaneous audience. The audience of sixty in A fits in the $238$ editions of the exact promise, and does not fit in the packing number $54$ at tolerance one. The audience of thirty in B fits at tolerance one, whose packing number is $32$, and not at tolerance two, whose number is $12$. A premiere to which everyone arrives at once can therefore promise separation by one slot in B and cannot in A, and the promise of Chapter~11 for locales has to be checked against the largest simultaneous audience in the same way.

\section{What a Repertoire Must Guarantee}\label{sec:guarantee13}

The preceding sections give three guarantees that a production can state about the editions that fill a person's life. The first is the \emph{novelty horizon}, the number of viewings of one production before a perceptual repeat, which is the packing number of the corrected pool and which \Cref{prop:life} converts into a catalogue. The second is \emph{simultaneous distinctness}, that no two viewers of a round hold equivalent editions, which is the Hall condition of Chapter~1 at the tolerance chosen and which \Cref{prop:cyclic} meets for the exact promise. The third is \emph{survival}, that these guarantees hold through the corrections of Chapters~7 to~10. The schedule is drawn from the pool that survives every anticipated amendment, the invariant pool of Chapter~9, or else repaired as in Chapter~10 when an amendment arrives, and the repair rewrites the ranks of the viewers it moves. The no-repeat guarantee of the cyclic schedule does not by itself survive a reallocation: a replacement that is valid today can duplicate an edition the viewer watched yesterday. For the lifetime guarantee to survive corrections, the allocation needs each viewer's playback history as a constraint, with the perceptual neighbors of past editions excluded when the promise uses a threshold.

The ledger of one viewer is the lifetime form of the record of Chapter~12. In the compact encoding it holds $17$ bits for each of the $4{,}301$ viewings of production A, which is $73{,}117$ bits or about nine kilobytes, and $14$ bits and about seven and a half kilobytes for B. At fourteen hours a week the same ledger has $30{,}112$ records and sixty-four kilobytes in A. These are the sizes at the premiere pool and would be smaller in the corrected pool. Production C has $28$ bits a record and about fifteen kilobytes for a life, for a production with no perceptual choice at all.

\section{What the Chapter Establishes and What It Assumes}\label{sec:limits13}

The propositions hold for any pool, any tolerance, and any viewing budget. The counts rest on assumptions stated beside them: that a viewing lasts as long as the measured passage, that the viewing rates are examples, that every production in the catalogue has the novelty horizon of A or of B, that the packing numbers of Chapter~2 on a Hamming dissimilarity stand in for measured ones, and that a revisit counts as new whenever the edition is perceptually new, which a viewer who remembers the story will not grant. The cyclic schedule treats the audience as one cohort viewing in rounds, and a real audience arrives and leaves.

The next chapter drops the assumption that the films of a repertoire are items. Chapter~14 asks what a later film does to an earlier one, so that releasing the same film in the tenth year of a viewer's life may not be the same as releasing it in the first.

The scripts \texttt{lawnmodel.py} and \texttt{verify\_ch13.py} reproduce each figure of this chapter.

\begin{thebibliography}{99}

\bibitem{flyxion2026gf}
Flyxion.
\newblock \emph{Generative Film: Categories, Sheaves, and Types for Cinema}.
\newblock Unpublished monograph, October 2026.

\end{thebibliography}

\end{document}
